\section{Ternary parent classes and marked neighbor trees}\label{sec:ternary-general}
Let $G$ be any positive definite symmetric rational ternary Gram
matrix, scaled to integral entries, and $\Gamma=\mathrm{SO}(L)$.
Proper and ordinary counts coincide in dimension three.
Theorem~\ref{th:incidence} gives the coefficients at every index;
at unramified odd primes its local traces can be organized on trees
whose vertices retain their global parent classes.

\paragraph{All local modules.}
Write $X_a=(1-aT)^{-1}$, $Y=(1-T^2)^{-1}$ and
$J_p(c)=X_1^3$ or $X_1Y$ according as $(c/p)=1$ or $-1$.
Let $u$ be a rotation preserving a lattice $E$. For $u\ne1$ let $U$
be the rational projection onto its positive axis. Put
$\delta=1$ if $U E_p\subset E_p$, and $\delta=0$ otherwise.
Write $S_p=(1+(p-1)T)X_1^2$ and $A_p=1+TS_p$.
The complete stable-subgroup table by order of $u$ is
\[
\begin{array}{c|l|l}
\operatorname{ord}u&\text{prime condition}&\sum_{k\ge0}b_{E,u}(p^k)T^k\\ \hline
1&\text{all}&X_1 X_p X_{p^2}\\
2&p\ne2& X_1^2 X_p\\
2&p=2,\ \delta=1&(1+3T)X_1^2 X_2\\
2&p=2,\ \delta=0&(1-T+4T^2)X_1^2 X_2\\
4&p\ne2& J_p(-1)\\
4&p=2& \delta S_2+(1-\delta)A_2\\
3&p>3& J_p(-3)\\
3&p=2&X_1Y\\
3&p=3&\delta S_3+(1-\delta)A_3\\
6&p>3&J_p(-3)\\
6&p=2&X_1Y\\
6&p=3&X_1^2 .
\end{array}
\]
Every rational proper torsion rotation in dimension three has one of
the listed orders: its quadratic complementary characteristic polynomial
is cyclotomic or a power of $X+1$. At the three exceptional
order/prime pairs $(2,2),(4,2),(3,3)$ the integral algebra generated
by $u$ is a fiber product over the residue field, and the split
condition is precisely integral preservation by $U$. The fiber-product calculation of Section~\ref{sec:cubic} applies to
an arbitrary stable $E_p$. Saturation under the two eigenprojections
gives the split module; the original module is either the split module
or a graph of residue components. In the half-turn case every graph of
a nonzero linear form on the two-dimensional negative residue plane
is conjugate under $\mathrm{GL}_2(\mathbb Z_2)$. In the Gaussian or
Eisenstein case multiplication by a unit accounts for each nonzero
residue scalar. Thus the counts are independent of any further
embedding data. The unramified and ramified direct-product rows
follow from the same Bezout splitting and the ideal series of the
degree-one or degree-two valuation rings; specifically at $(6,3)$
the factors $X-1,\Phi_6(X)$ are relatively prime over $\mathbb Z_3$.
The projections in the three exceptional cases are
$(1+u)/2$, $(1+u^2)/2$, and $(1+u+u^2)/3$, respectively. These calculations give the local formula at all primes, including
those dividing $\det G$.

Enumerate parent transporters by integral matrices $B$ satisfying
$B^tGB=m^2G$ and $\det B=m^3$, subject to the column bounds
$|B_{ij}|\le m\sqrt{G_{jj}(G^{-1})_{ii}}$, from Cauchy--Schwarz.
These matrices give the parent transporters $B/m$. Identify parents
by right multiplication by $\Gamma$ and retain only
$[L:L\cap (B/m)L]\mid m$. Each retained parent contains $mL$ and
lies in $m^{-1}L$. For a fixed torsion transporter $u$, any preserved
parent subset is a union of cycles entirely inside the finite retained
list; it must include the cycle of $L$. Successively adding or omitting
the other cycles gives the sums of Theorem~\ref{th:incidence}. The projection
test in the table uses a local lattice basis, or reduction of the
intersection modulo a sufficiently large $p$-power. 

\paragraph{Prime powers and global parent labels.}
Take $p\nmid 2\det G$. In the self-dual neighbor tree at $p$, identify
a vertex with the global lattice given by it at $p$ and $L$ at all
other places. Mark a vertex if this global lattice is properly
isometric to $L$. For a finite nonempty subtree $\mathcal T$, write $\partial\mathcal T$ for its set
of leaves (including the unique vertex of a singleton), and let
$E_{\mathcal T}$ be the global intersection of its vertices. Write
$p^{j_{\mathcal T}}=[L:E_{\mathcal T}]$ for trees containing $L$.
\begin{theorem}\label{th:ternary-hull}
For every $k\ge0$ the actual coefficient at $p^k$ is
\[
 \frac1{|\Gamma|}
 \sum_{\mathcal T}\frac1{|\partial\mathcal T|}
  \sum_{u\in\mathrm{Stab}^+(\mathcal T)}
      (-1)^{c(u,\partial\mathcal T)-1}
      b_{E_{\mathcal T},u}\big(p^{\,k-j_{\mathcal T}}\big).
\]
Sum over subtrees with $L$ a leaf, $j_{\mathcal T}\le k$, every leaf marked, and
no marked interior vertex of the chosen subtree. The stabilizer consists of global
proper rational isometries preserving $\mathcal T$ and fixing $L_\ell$
at all places $\ell\ne p$; $c$ counts cycles on leaves. This is a
finite sum, with $b$ given by the table above. If the genus has
one class, only singletons and edges survive.
\end{theorem}
\begin{proof}
In the unanchored incidence sum proved in Theorem~\ref{th:incidence},
take the convex hull of each parent set $F$ in the $p$-tree. Its
leaves must lie in $F$. Conversely every subset of its marked
vertices containing its leaves spans the tree. Intersecting all the
vertices of the hull gives exactly the intersection of the original
subset, since the unique geodesic between two vertices contains
their intersection in every one of its lattices. Every hull vertex is
on such a geodesic or is itself a leaf.

Within a fixed hull, orbit--stabilizer rewrites the unanchored
sum over orbits of $F$,
$\sum_{[F]}|\mathrm{Stab}(F)|^{-1}\sum_{u\in\mathrm{Stab}(F)}$
as $|\mathrm{Stab}(\mathcal T)|^{-1}\sum_{u\in\mathrm{Stab}(\mathcal T)}
\sum_{F:uF=F}$, with the same integrands. For each $u$ the optional
marked interior vertices form a union of cycles. Including an
optional cycle changes the sign; the sum therefore vanishes if any
such vertex exists. Otherwise only the leaf set remains. Marking a
leaf as $L$ replaces the unanchored stabilizer denominator by
$|\partial\mathcal T|\,|\Gamma|$ by the same orbit-stabilizer identity. The
restriction $j_{\mathcal T}\le k$ bounds every vertex's distance to $L$ by
$k$, by the symmetric Smith indices of the self-dual tree, so the
tree sum is finite.
\end{proof}

The leaf sets also admit a graph-theoretic description. In the ball of
radius $k$ around $L$, join two marked vertices $P,Q$ in an auxiliary
graph if there is no third marked vertex $R$ satisfying
$\mathrm{dist}(P,R)+\mathrm{dist}(R,Q)=\mathrm{dist}(P,Q)$.
The leaf sets in the formula are exactly the cliques containing $L$
whose intersection index divides $p^k$. Indeed the ball is convex,
every interior vertex lies on a path between two leaves, and any
member of a proposed leaf set that were interior would lie between
two other members. The distances are obtained from the denominator of
the relative transporter in a basis of $L$ by the self-dual Smith
exponents. The marked vertices therefore determine the leaf sets.

At $k=1$ the formula reduces to a class-sensitive expression. Let
$A_\Gamma(p)$ count the $\Gamma$-orbits of hyperplanes of $L/pL$.
Let $t_p$ be the number of double cosets $\Gamma g\Gamma$ of
rational proper isometries with $[L:L\cap gL]=p$, and $r_p$ the
number of these cosets fixed by inversion. Then
\[
                a_L(p)=A_\Gamma(p)-\tfrac12(t_p-r_p).
\]
Indeed the double cosets are the orbits of marked outgoing neighbor
edges. Reversal induces inversion of double cosets. The intersection
child of a fixed orbit contributes one ambient orbit and one true
class. Each pair of distinct reversed orbits contributes two ambient
orbits and one true class. All other hyperplane children have just
one marked parent. Even two distinct incident neighbors have distinct
intersection hyperplanes by the local tree construction. For any
$p$ here, including $3$, the first term is the average over $u\in\Gamma$
of $\sum_{\alpha\in\mathbb F_p^\times}
(p^{\dim\ker(u-\alpha)}-1)/(p-1)$, since it counts invariant
hyperplanes via eigenlines in the dual. For $p>3$ the contributions
are $p^2+p+1$, $p+2$, $2+(-1/p)$ and $2+(-3/p)$ for orders
$1$, $2$, $4$ and $3$ or $6$, respectively.

For indices supported on several good primes the same cancellation
has an exact finite form. Exclude $m=1$, and for each $p\mid m$ take a
finite subtree $\mathcal T_p$. Put $\mathcal Q=\prod_{p\mid m}\mathcal T_p$, with $L$ prescribed
away from these primes. Mark a tuple if its global lattice is properly
isometric to $L$. For $u$ preserving the box and the external local
lattices the unanchored weight is
\[
 \sum_{\mathcal U}(-1)^{|\mathcal U|-1}.
\]
Here $\mathcal U$ runs over collections of cycles of $u$ on marked
tuples, whose projections together contain every leaf of every factor
(a singleton counts as a leaf). The whole contribution is this weight
times $b_{E_{\mathcal Q},u}(m/[P:E_{\mathcal Q}])$, summed over $u$ and divided by the
global proper stabilizer order of $\mathcal Q$; $P$ is any vertex and $E_{\mathcal Q}$ the
intersection of all vertices. Use one representative of each global
box class that contains a marked tuple, with $[P:E_{\mathcal Q}]\mid m$.
Indeed coordinatewise convex hull gives the same intersection at
each prime; a union of cycles spans the box precisely by this leaf
test. Marking one tuple bounds all relevant trees by the corresponding
index exponents and hence proves finiteness, as above. The parent
class marks occur on the global tuples, and need not factor as
independent choices in the local trees. At the remaining primes the
all-index incidence expression and the explicit stable-module table
apply without any tree assumption.

For $L=\langle1,1,11\rangle$ and $p=3$, the line spanned by $(1,0,1)$
modulo three is isotropic. The corresponding neighbor cannot be
isometric to $L$: its intersection hyperplane omits $e_1$, and the
only rational unit vectors with denominator dividing three are
$\pm e_1,\pm e_2$ (multiply the norm equation by nine). Thus it
has at most two unit vectors whereas $L$ has four. The local
neighbor is nevertheless in the proper genus. Consequently it is an
unmarked vertex in the theorem.

For $G=\operatorname{diag}(1,1,7)$ and $p=3$, a marked vertex
can occur beyond an unmarked intermediate vertex. Every rational
unit vector with denominator dividing three is one of $\pm e_1,\pm e_2$ or
$(\pm1,\pm1,\pm1)/3$. No vector of the latter type is perpendicular
to either unit coordinate vector or any other vector of the latter
type. Thus a proper isometry of denominator dividing three must send
$e_1,e_2$ to the coordinate vectors and hence be integral. No adjacent
vertex is marked. But
\[
 A=\begin{pmatrix}-7&-2&-14\\-2&-7&14\\-2&2&5\end{pmatrix},
 \qquad A^tGA=81G,\qquad\det A=729 .
\]
The rational rotation $A/9$ has exact denominator nine, so its marked
vertex lies at distance two. The path between these two marked
endpoints has precisely an unmarked interior vertex.
